\ficha{Marinho (2021), Metalistas e papelistas na imprensa}{marinho2021}

\begin{identificacao}
MARINHO, Giullia Bueno. \emph{``Oh! Quantas species!...'': ``metalistas'' e
``papelistas'' na imprensa brasileira, 1888-1892}. 2021. Monografia
(Bacharelado em Ciências Econômicas) - Faculdade de Economia, Administração e
Contabilidade, Universidade de São Paulo, São Paulo, 2021. Orientador: José
Flávio Motta.

Monografia de graduação, 96 páginas, em português. Leitura seletiva: capítulos
1, 2, 3, 6 e 7 por inteiro; capítulos 4 e 5, os empíricos, por amostragem.
\end{identificacao}

\bloco{Problema e tese}
A pergunta é se o debate entre metalistas e papelistas, tal como a
historiografia o define, aparece na imprensa brasileira do quinquênio
1888-1892, e se as posições dos jornais podem ser identificadas com interesses
de grupos ou de regiões. A resposta tem duas partes. O debate aparece, mas
raramente em forma explícita, e a leitura precisa ser, nas palavras da autora,
uma avaliação ``não binária'' (p. 71). E a associação entre posição e interesse
de classe, proposta por Fonseca e Mollo (2012), não se confirma no material:
os órgãos mais ligados ao comércio e à produção defenderam a conversibilidade
metálica, não a emissão.

\bloco{Argumento}
A monografia propõe uma cronologia das reformas bancárias de 1888 a 1892,
acompanhada, decreto a decreto, pela repercussão em pouco mais de uma dezena
de periódicos de províncias e estados distintos, e encerra com uma
classificação de cada jornal nos dois eixos do debate. O que ela quer mostrar
é que a classificação é possível, mas custosa, e que a hipótese sociológica
corrente sobre quem defendia o quê falha no teste.

\begin{enumerate}[leftmargin=*,nosep]
\item A genealogia é a de Gremaud e a de Fonseca e Mollo, reproduzida sem
acréscimo: bulionistas e antibulionistas, Currency School e Banking School, e o
padrão-ouro como vitória institucional da ortodoxia, o que explica por que o
metalismo se tornou enraizado e de difícil superação (p. 10-11). A adesão
social esperada também vem deles: rentistas e proprietários de terra do lado
metalista, industriais, comerciantes e produtores rurais do lado papelista
(p. 11, p. 12).

\item A digressão metodológica é curta e defensiva. Os acervos são a
Hemeroteca Digital da Biblioteca Nacional e o do Estado de S. Paulo (p. 20).
A autora invoca Capelato, Lapuente e Mariani para reconhecer que o periódico
está ``envolvido em um jogo de interesses'' (p. 20) e que a página constrói
consensos ilusórios (p. 21), e declara que buscará atos oficiais, editoriais
com posicionamento dos veículos e opiniões anônimas e influentes (p. 25). Não
há protocolo de busca, termo de pesquisa, recorte de datas, contagem de peças
nem critério declarado de seleção da peça individual.

\item Os capítulos 4 e 5 são narrativos: Auxílios à Lavoura e Lei Bancária de
24 de novembro de 1888 sob João Alfredo, a supressão do lastro em apólices por
Ouro Preto, o decreto de 17 de janeiro de 1890 de Rui Barbosa, com lastro em
títulos públicos e bancos regionais, e os recuos sucessivos até o Banco da
República e o Encilhamento (p. 86-87). A tese política que os atravessa é a
inconstância: os ministros inclinavam-se ao metalismo ou ao papelismo, ao
monopólio ou à pluralidade, ``conforme lhe fosse mais conveniente'' (p. 86).

\item O capítulo 6 classifica os jornais em dois eixos cruzados,
conversibilidade e faculdade emissora, e o cruzamento é o achado formal do
trabalho. Metalista com monopólio: Jornal do Commercio. Metalista com
pluralidade: A Nação e Correio Paulistano. Papelista com pluralidade: Diário
de Notícias e A Província de São Paulo. Papelista com monopólio: O Paiz.
Papelista com posição ambígua sobre a emissão: Gazeta de Notícias e O Cruzeiro
(Tabela 3, p. 83). Os oito da tabela não esgotam os jornais usados.

\item O teste da hipótese de interesse falha. O representante da Associação
Comercial do Rio de Janeiro pede na Gazeta de Notícias a volta da moeda
conversível, e o Jornal do Commercio, órgão do comércio por excelência, é o
metalista mais firme e mais consistente do corpus (p. 72, p. 74). A explicação
oferecida é cultural, não econômica: o ``espectro metalista'', a dificuldade de
romper com um sistema tido como pilar da civilização (p. 72-73).
\end{enumerate}

\bloco{Conceitos e definições operacionais}
Metalismo: conversibilidade da moeda a uma quantidade fixa de metal, com
estabilidade cambial como objetivo e o Estado obrigado a manter a oferta
monetária equivalente às reservas metálicas, impedindo o ágio sobre o ouro
(p. 84). Papelismo: oferta monetária como instrumento de crescimento, câmbio
como variável consequência e não de ajuste, juros como o verdadeiro barômetro
(p. 84). Faculdade emissora: eixo separado do primeiro, monopólio contra
pluralidade bancária, e é ele que produz os cruzamentos inesperados da Tabela
3 (p. 83). Espectro metalista: nome que a autora dá à dificuldade de abandonar
o conservadorismo econômico (p. 73, nota 25), e que funciona como explicação
residual sempre que o interesse material e a posição publicada divergem.
Avaliação não binária: reconhecimento de que a posição do jornal se infere de
uma sequência de colunas e seções, não de uma declaração (p. 71).

\bloco{Evidência e método}
Fontes primárias: periódicos do Rio de Janeiro (Jornal do Commercio, Gazeta de
Notícias, O Paiz, Diário de Notícias, O Cruzeiro, A Nação, Gazeta da Tarde),
de São Paulo (A Província de São Paulo, Correio Paulistano, Diário Mercantil),
de Pernambuco (Diário de Pernambuco, Jornal do Recife) e do Rio Grande do
Norte (O Povo), mais os Retrospectos Commerciaes do Jornal do Commercio, que
são a espinha dorsal factual dos capítulos 4 e 5 (p. 74). O critério de
seleção dos jornais é misto e declarado: importância e circulação para os
grandes diários, contraposição deliberada de pares rivais (A Nação contra O
Cruzeiro, Jornal do Recife contra Diário de Pernambuco) e cobertura regional
do Nordeste, esta última com expectativa anunciada de antemão, ``em um
prognóstico imaturo, opiniões mais conservadoras de tradicionais proprietários
de terra'' (p. 24). Não há amostragem, censo de edições, série temporal nem
contagem. O capítulo empírico se intitula ``uma breve análise'' e a evidência
é um conjunto de citações escolhidas, cada uma com data e quase sempre com
página.

A distinção entre posição do veículo e matéria hospedada é enunciada no
capítulo 3 e não é sustentada na prática. A peça do Diário de Pernambuco de 25
de julho de 1888, que funda a leitura papelista do jornal, saiu em
``Publicações a pedido'', espaço de terceiros (p. 27-28). A citação que a
Tabela 3 dá como exemplo da Gazeta de Notícias, de 15 de junho de 1888, é
transcrição de discurso do conselheiro Matta Machado em sessão parlamentar
(p. 28, p. 83). Várias passagens atribuídas ao Jornal do Commercio são
transcrições da Cidade do Rio ou do Diário Mercantil, e uma delas é lida na
Gazeta da Tarde (p. 75-76, p. 82, notas 32, 35 e 56). A autora registra que O
Paiz de 18 de janeiro de 1892 traz, na mesma edição, um editorial transcrito
de O Tempo que ataca Rui Barbosa e uma coluna que o exalta (p. 47), e ainda
assim classifica O Paiz como papelista com base na segunda.

Quanto aos interesses, as associações propostas são: Jornal do Commercio,
conservadorismo doutrinário e público comercial (p. 74); Diário de Notícias e
O Paiz, proximidade pessoal com Rui Barbosa (p. 77, p. 80); Diário de
Pernambuco, reivindicação de recursos para um Nordeste decadente, contra o
Jornal do Recife, ligado à antiga aristocracia (p. 78-79, p. 89); Correio
Paulistano e Diário Mercantil, cafeicultores e comerciantes paulistas exigindo
que a emissão chegasse a São Paulo, pressão atendida pelo decreto de 31 de
janeiro de 1890, que criou uma quarta região emissora (p. 81-82).

\bloco{Agentes e vertentes}
Metalistas: Visconde de Ouro Preto, Torres Homem, Francisco Belisário Soares
de Souza e Joaquim Murtinho (p. 11). Papelistas: Mauá, Souza Franco e Rui
Barbosa (p. 12). Financistas acusados de capturar as reformas: Conde de
Figueiredo, junto a Ouro Preto, e Francisco de Paula Mayrink, junto a Rui
Barbosa (p. 86). Campos Sales lidera a pressão paulista por uma região
emissora própria (p. 82). Literatura mobilizada: Fonseca e Mollo (2012),
Gremaud (1997), Franco (2005), Gambi (2015), Salomão (2017), Villela (2001) e
Suzigan (1986).

\bloco{Críticas e limites}
O trabalho não mede nada. Não há denominador, não se sabe quantas edições
foram vistas nem quantas peças sustentam cada rótulo, e a Tabela 3 sintetiza
oito jornais quando o texto usa mais de dez. A tabela contradiz a si mesma no
caso da Gazeta de Notícias, cuja descrição a diz ``defensora do monopólio
emissor'' enquanto a coluna própria registra ``ambíguo'' (p. 83), o que o
corpo do capítulo confirma ao mostrar manifestações nos dois sentidos (p. 78).
A atribuição de posição por proximidade pessoal, com O Paiz e o Diário de
Notícias classificados como papelistas em boa parte porque Rui Barbosa
escreveu neles, confunde alinhamento com um agente e posição sobre o objeto
monetário (p. 77, p. 80). O espectro metalista opera como explicação não
falseável do resultado negativo do teste de interesse. Nada disso invalida o
achado principal, que é justamente negativo e vale como alerta: no material de
imprensa, a correspondência entre classe e doutrina não se verifica.

\chave{O que se observou foi uma dificuldade de encontrar posicionamentos
explícitos em defesa de uma ou outra corrente do debate estudado. Os jornais
publicaram com abundância, de maneira informativa, as modificações e decretos
da esfera econômica, no entanto, os apoios, críticas e opiniões a eles foram
mais difíceis de serem encontrados}{p. 88}

\begin{dialogo}
\item Reaproveitar o eixo duplo. Marinho separa a posição sobre a
conversibilidade da posição sobre quem emite, e é o segundo eixo que produz
suas combinações inesperadas (p. 83). O análogo em 1906 a 1914 é separar a
posição sobre a Caixa de Conversão da posição sobre a competência dela frente
ao Tesouro e ao Banco do Brasil e sobre o limite de emissão. Testa se um mesmo
jornal pode ser favorável à Caixa e contrário à ampliação do limite, o que
derrubaria qualquer escala unidimensional de posição.

\item Testar de novo a hipótese de interesse, com denominador. Marinho a
refuta com citações escolhidas. Verificar no corpus se o Correio Paulistano,
órgão do PRP e dos cafeicultores, defende a taxa baixa de Taubaté por
interesse exportador, e se O Paiz e a Gazeta de Notícias, cariocas e ligados à
praça, divergem dele. Se a clivagem por região ou por interesse aparecer com
contagem, o resultado contraria Marinho e devolve fôlego a Fonseca e Mollo; se
não aparecer, reforça o achado dela por método melhor.

\item Vigilância sobre voz, com o caso dela como corner case a evitar. A peça
do Diário de Pernambuco vem de ``Publicações a pedido'' e o exemplo da Gazeta
de Notícias é discurso parlamentar transcrito (p. 28, p. 83). É exatamente o
erro que o campo de voz do pipeline existe para impedir. Procurar em 1906 a
1914 a mesma armadilha na SECÇÃO LIVRE e nos relatos de sessão, e reportar que
fração das peças hospedadas seria classificada como posição do jornal caso a
voz não fosse marcada.

\item Reprodução entre jornais. Marinho cita o Jornal do Commercio pela Gazeta
da Tarde e o Diário Mercantil pelo Jornal do Commercio (notas 32 e 56). É
confirmação, em fonte independente, de que a mesma peça circula entre diários
no período. Medir a taxa de reprodução no corpus de 1906 antes de tratar cada
ocorrência como opinião distinta de um jornal.

\item Sobrevida do vocabulário. Marinho mostra que em 1888-1892 os próprios
jornais quase nunca se dizem metalistas ou papelistas, e que a rotulação é
sempre do analista (p. 88). Verificar se em 1906 a 1914 os termos aparecem em
autodescrição ou apenas como acusação. Um resultado igual ao dela, rótulo do
analista e não da fonte, reforça a decisão da dissertação de tratar o par como
vocabulário de época e usar o eixo ortodoxo contra expansionista como síntese
analítica.
\end{dialogo}

\usonocapitulo{Parte III, ao lado de Fonseca e Mollo (2012), e sobretudo como
ponte metodológica. É o trabalho mais próximo do desenho da dissertação: mesma
fonte, a imprensa; mesmo par de categorias; mesmo tipo de pergunta sobre a
associação entre posição publicada e interesse. Serve a três usos. Primeiro,
como precedente de que a leitura de posição em jornal do período é viável e de
que a classificação binária não é, o que ampara a decisão de tratar metalismo
e papelismo como vocabulário de época. Segundo, como evidência contrária à
hipótese sociológica de Fonseca e Mollo, útil para o capítulo afirmar que a
correspondência entre classe e doutrina é conjectura ainda não verificada em
imprensa. Terceiro, e principalmente, como contraste de método: o que a
dissertação acrescenta ao que já existe é o denominador, a unidade de análise
declarada, a distinção sistemática de voz e a conferência mecânica de citação,
tudo ausente aqui. O trabalho para em 1892 e não trata da Caixa de Conversão;
nenhuma afirmação sobre 1906 a 1914 pode ser apoiada nele.}
